\documentclass[11pt,letterpaper]{article}
\usepackage[margin=1in]{geometry}
\usepackage[T1]{fontenc}
\usepackage{lmodern}
\usepackage{amsmath,amssymb}
\usepackage{xcolor}
\usepackage{enumitem}
\usepackage{titlesec}
\usepackage{hyperref}
% Article handout: reuse SDSU colors without loading the Beamer-only theme.
\definecolor{SDSUDarkRed}{HTML}{A6192E}
\hypersetup{colorlinks=true,urlcolor=SDSUDarkRed,linkcolor=SDSUDarkRed,
  pdftitle={AE 696 Project Proposal Template},pdfauthor={Dr. Jun Chen}}
\setlength{\parindent}{0pt}
\setlength{\parskip}{4pt}
\setlist[itemize]{leftmargin=1.35em,itemsep=2pt,topsep=3pt,parsep=0pt}
\titleformat{\section}{\normalsize\bfseries}{}{0pt}{}
\titlespacing*{\section}{0pt}{9pt}{3pt}
\newcommand{\prompt}[1]{{\itshape #1}}
\begin{document}

\begin{center}
{\large\bfseries\color{SDSUDarkRed} AE 696 Project Proposal}\\[3pt]
Reinforcement Learning with AI Applications in Engineering\\
Dr. Jun Chen \quad | \quad Fall 2026
\end{center}

\textbf{Purpose.} Propose a feasible engineering application of reinforcement
learning. Define the problem, a preliminary RL formulation, and a measurable
experimental plan. No code or preliminary results are required at this stage.

\textbf{Submission.} Work individually or in a team of two. Submit one PDF per
project on Canvas by the posted Week 9 deadline. Use 1 to 2 pages for the proposal
body, excluding references, with 11 pt or larger text and 1 inch margins.

\textbf{Using this template.} Replace the italic prompts with your responses;
delete these opening instructions from your submission. Keep the numbered
headings. Short paragraphs or bullet points are acceptable. Algorithm and reward
choices may be preliminary, but must be specific enough to assess feasibility.

\textbf{Project title:} \prompt{[Title]}\\
\textbf{Student name(s) and SDSU email(s):} \prompt{[Name and email for each member]}

\section*{1. Engineering Problem and Objective}
\prompt{Describe the system, the task, and the engineering objective in a short
paragraph. Explain why learning through sequential interaction is useful here.
Identify what you will contribute, such as a task modification, reward design,
controller comparison, or robustness study. An existing simulator or algorithm
is acceptable; merely reproducing a tutorial is not a sufficient project.}

\section*{2. Preliminary MDP or POMDP Formulation}
\prompt{Specify each item below. Give variables, dimensions, units, or bounds
where applicable; identify any details that remain to be determined.}
\begin{itemize}
  \item \textbf{State and observation:} \prompt{What describes the system, and
  what information does the agent actually receive? If partially observable,
  state whether you will use history, memory, or state estimation.}
  \item \textbf{Action:} \prompt{What does the agent command? Specify discrete
  choices or continuous dimensions and bounds.}
  \item \textbf{Reward:} \prompt{Give a preliminary equation or an explicit list
  of reward terms, signs, and intended effects. List safety constraints separately.}
  \item \textbf{Dynamics and episodes:} \prompt{Identify the simulator or transition
  process and whether its model is available to the algorithm. Specify initial
  conditions, success/failure termination, time limit, and discount factor.}
\end{itemize}

\section*{3. Method and Implementation Plan}
\prompt{Name the initial RL algorithm and explain its suitability for the action
space and task. Name at least one comparison baseline, such as a random policy,
heuristic, classical controller, or another RL method. Identify the environment,
software libraries, available computing resources, and any data needed. Distinguish
existing code you will reuse from the components you will develop or modify.}

\section*{4. Evaluation Plan}
\prompt{Describe the planned experiments, not predicted numerical outcomes.}
\begin{itemize}
  \item \textbf{Metrics:} \prompt{Specify at least two quantitative metrics,
  including one tied to the engineering objective, such as success rate, tracking
  error, control effort, or constraint violations. Define how each is measured.}
  \item \textbf{Protocol:} \prompt{Specify the training budget, number of independent
  training seeds, and evaluation episodes per trained policy. A suggested starting
  point is 3 training seeds and 20 evaluation episodes per seed; explain a smaller
  plan if computing limits require it. Evaluate frozen policies without learning.}
  \item \textbf{Comparison and outputs:} \prompt{Use the same test conditions and
  metrics for the RL policy and baseline. State how results will be summarized,
  including variability across training runs. Name the planned plots or tables
  and a measurable target or hypothesis to test; improved performance is not
  guaranteed or required.}
\end{itemize}

\section*{5. Milestones, Responsibilities, and Backup Plan}
\prompt{Give target weeks or dates for environment setup and baseline testing,
initial RL training, main experiments, and final analysis. Use the Canvas schedule
for the progress check-in, presentation, and report deadlines. For teams, assign
specific responsibilities to each member. Identify the main technical risk and
a simpler backup that preserves a meaningful RL experiment.}

\section*{6. References and Disclosure}
\prompt{Cite the papers, software, environments, and tutorials you plan to use.
Give enough information or links to locate them. Disclose any substantial
AI-generated proposal text or code, as required by the course policy. All members
must understand and be able to explain the submitted work.}

\section*{Submission Check}
\prompt{Delete this checklist before submission. Confirm that all six sections
are completed; states, observations, actions, rewards, and episode definitions are
specified; a baseline and at least two metrics are named; the experiment budget,
milestones, responsibilities, and backup plan are stated; and sources and any
substantial AI assistance are acknowledged. Check the PDF for readability and
the proposal body for the page limit.}

\textbf{Instructor review.} Proposals will be reviewed for completeness,
technical consistency, feasible scope, and a testable evaluation plan.
Instructor feedback may require revisions before the project scope is approved.

\end{document}
